\exerciseeditionsection{Solutions to Wald's Problems}
\ExerciseEditorialNote

\begin{enumerate}[leftmargin=2.8em,itemsep=1.2em]
\WaldSolution{1}
For the Killing field \(\psi^a\),
\[
 \mathcal L_\psi\epsilon_{abcd}=(\nabla_e\psi^e)\epsilon_{abcd}=0.
\]
The pullback by an isometry preserves the Levi--Civita derivative, so its Lie
derivative commutes with \(\nabla_a\).  Apply \(\mathcal L_\psi\) to
\(\omega_a=\epsilon_{abcd}\xi^b\nabla^c\xi^d\).  The volume-form term
vanishes as above, and every remaining term contains
\(\mathcal L_\psi\xi^a=0\) or its derivative.  Therefore
\(\mathcal L_\psi\omega_a=0\).

\WaldSolution{2}
For the Papapetrou metric, the determinant and inverse of the Killing-orbit
block are
\[
 \sqrt{-g}=\rho V^{-1}e^{2\gamma},qquad
 (g^{AB})_{A,B=t,\phi}
 =\begin{pmatrix}
 -V^{-1}+Vw^2/\rho^2&Vw/\rho^2\\
 Vw/\rho^2&V/\rho^2
 \end{pmatrix}.
\]
Substitution in the coordinate Ricci formula separates the vacuum equations
into two second-order equations for \(V,w\) and two first-order equations for
\(\gamma\):
\[
 \boldsymbol\nabla\mathbin\cdot
 \left(V^{-1}\boldsymbol\nabla V+
ho^{-2}V^2w\boldsymbol\nabla w\right)=0,
 \qquad
 \boldsymbol\nabla\mathbin\cdot
 \left(\rho^{-2}V^2\boldsymbol\nabla w\right)=0,
\]
\[
 \partial_\rho\gamma
 =\frac{\rho}{4V^2}(V_{,\rho}^2-V_{,z}^2)
 -\frac{V^2}{4\rho}(w_{,\rho}^2-w_{,z}^2),
\]
\[
 \partial_z\gamma
 =\frac{\rho}{2V^2}V_{,\rho}V_{,z}
 -\frac{V^2}{2\rho}w_{,\rho}w_{,z}.
\]
These are Eqs.~\eqref{eq:7.1.24}--\eqref{eq:7.1.27}.  Differentiating the
last pair in opposite orders shows that their integrability condition is the
first pair.

\WaldSolution{3}
\begin{enumerate}[label=\alph*),leftmargin=2em]
\item Since \(u_a=N\nabla_at\) and the axial Killing field is tangent to every
\(t=\text{constant}\) hypersurface,
\[
 L=u^a\psi_a=N\psi^a\nabla_at=0.
\]

\item Write \(u^a=u^t(\xi^a+\Omega\psi^a)\).  Then
\[
 0=u^a\psi_a=u^t(W+\Omega X),
\]
so
\[
 \Omega=\frac{\dd\phi}{\dd t}=-\frac{W}{X}.
\]
\end{enumerate}

\WaldSolution{4}
\begin{enumerate}[label=\alph*),leftmargin=2em]
\item At a point away from the center, the connected isotropy subgroup of
spherical symmetry consists of rotations about the radial direction.  The
only null directions fixed by all these rotations are the ingoing and
outgoing radial directions.  Every principal null direction must therefore
be one of these two, so at least one is repeated and the spacetime is
algebraically special.  In a static spacetime, time reflection interchanges
the two radial directions and preserves their multiplicities.  Four
principal directions are thus split equally between them: both are double,
which is type II--II.

\item At any event of a Robertson--Walker spacetime the connected isotropy
group is the full spatial rotation group.  No null direction is fixed by
every spatial rotation.  A nonzero Weyl tensor would supply a finite,
isometry-invariant set of principal null directions, each fixed by the
connected isotropy group.  This is impossible, so \(C_{abcd}=0\).
\end{enumerate}

\WaldSolution{5}
Let
\(\lambda=\xi^a\xi_a\) and
\(\omega_a=\nabla_a\omega=epsilon_{abcd}\xi^b\nabla^c\xi^d\).
The Killing identity in vacuum gives
\[
 \nabla_a\nabla_b\xi_c=R_{cbad}\xi^d,
 \qquad \nabla^a\nabla_a\xi_b=0.
\]
Differentiate
\[
 H_{ab}=2\lambda\nabla_a\xi_b
 +\omega\epsilon_{abcd}\nabla^c\xi^d
\]
and antisymmetrize over three indices.  Terms with second derivatives of
\(\xi\) reduce by the Killing identity and cancel by the vacuum Bianchi
identity.  In the remaining terms use
\(\nabla_a\lambda=2\xi^b\nabla_a\xi_b\),
\(\nabla_a\omega=\epsilon_{abcd}\xi^b\nabla^c\xi^d\), and the
epsilon-contraction identity.  The two quadratic products of
\(\nabla\xi\) cancel, leaving
\[
 \nabla_{[e}H_{ab]}=0,
\]
which is Eq.~\eqref{eq:7.4.8}.
\end{enumerate}
